% NMI Introduction: no \section heading (journal style).
Peer review is the gate through which scientific work passes, and it is under
strain: submissions at major venues have outgrown the reviewer pool, and the
process is demonstrably noisy. The NeurIPS consistency experiments found that a
large fraction of accept/reject decisions would flip under a re-run of the same
process, with the asymmetry that review is comparatively reliable at \emph{rejecting
weak work} but poor at \emph{ranking strong work}
\citep{cortes2021inconsistency,beygelzimer2023arbitrary}. Against this backdrop,
large language models are increasingly used to assist review---and publisher policy
is caught between uncontrolled, undisclosed use, which is risky, and categorical
prohibition, which forgoes potentially useful signal.

Most studies of AI in peer review evaluate the \emph{text} a model produces. A
large randomized study at ICLR~2025 showed that LLM feedback can make reviews more
informative and increase reviewer--author engagement \citep{thakkar2026feedback};
other work compares machine and human review comments for overlap, specificity, and
factuality \citep{liang2024llm,darcy2024marg,du2024llmsassist}. A parallel line
documents the risks of uncontrolled use: AI-assisted reviews already inflate scores
and acceptance rates at scale \citep{russolatona2024lottery}, and adversarial hidden
prompts embedded in manuscripts can manipulate AI-assisted review outright
\citep{lin2025hidden}. Almost none of this work asks the question deployment turns
on: whether the numeric \emph{score} such a system assigns is a valid signal of the
quality humans ultimately perceive. A score is what triages a queue or flags a
submission for a second look; used that way it must be validated as a measurement
against an external outcome, with the discipline a noisy ground truth demands.

We study \textsc{AIPR}, a system that reads a submitted PDF and produces five
0--100 quality dimensions and a weighted overall score, and we use it as the
instrument to test that question on public peer-review outcomes. Our ground truth
is the public record of a major machine-learning venue on OpenReview: the decision
tier of each submission (reject, poster, oral) and the mean numeric rating its
reviewers assigned. Because this venue publishes \emph{rejected} submissions, the
class our central claim concerns is fully observable. AIPR is deployed as an
interactive assistant; here we validate the baseline score it emits on first read,
before any interaction---the hardest case for an automated signal and the one a
triage use relies on.

Our claim is deliberately narrow, and complementary to the feedback-quality work
above: a feedback assistant can be useful even if it does not rank manuscripts,
whereas a triage score must be validated against the venue's own outcomes and bounded
to human oversight. The AIPR score (i) flags weak submissions---papers it scores low
are rejected far above the base rate---and (ii) shows strong agreement with human
outcomes across both the decision tiers and the continuous reviewer rating. We do
\emph{not} claim it predicts the acceptance decision, ranks strong papers, or
substitutes for a reviewer; the bound mirrors the ``reliable at the bottom, noisy at
the top'' structure of human review itself.

A validated score raises an obvious question: a frontier model is available to
anyone, so what does the engineered pipeline add over simply asking that model? We
test AIPR against a generic one-paragraph prompt on the identical model and PDF. The
bare prompt, with no rubric and no audit, already separates rejected from accepted
submissions on its own; the full pipeline's discrimination is higher, but the planned
paired comparison is not statistically resolved. The lesson is about the model, not
the pipeline---with or without our prompt, a frontier model already produces a
first-pass reviewing signal that substantially tracks the human outcome in this
setting, so raw model intelligence is not the observed bottleneck, and the reviewer
keeps the decision. What the engineering adds is therefore not validity but
\emph{deployability}: a score reliable across repeated runs that arrives with a
grounded, auditable, rubric-anchored review. The right policy question is then not
whether AI can ever assist review, but which controlled, disclosed, auditable,
human-accountable systems should.

The study's standard of evidence is set by three practices: the hypotheses, primary
metric, and low-score threshold were pre-registered before any score was joined to
any outcome; every number is generated from the released data by a single script;
and we report the score's failure modes---the rejected papers it rated highly and the
unconventional accepted papers it underrated---rather than headlining best-case
agreement.
